\documentclass[10pt,a4paper]{amsart}
\usepackage[margin=19mm]{geometry}
\usepackage{amsmath,amssymb,mathtools}
\usepackage[scr=rsfs,cal=euler]{mathalfa}
\usepackage{xfrac}
\usepackage[unicode,hidelinks,bookmarks=false]{hyperref}
\newcommand{\ie}{i.e.\ }
\newcommand{\cf}{cf.\ }
\newcommand{\resp}{respectively\ }
\newcommand{\Subsection}{\subsection}
\providecommand{\nobreakdash}{\nobreak\hbox{-}}
\setlength{\emergencystretch}{2em}
\multlinegap0pt
\usepackage{bm}
\usepackage{bbm}
\usepackage{dsfont}
\usepackage{xspace}
\usepackage{cancel}
\usepackage{mathtools}
\usepackage{amsthm}
\makeatletter
\@ifundefined{definition}{\newtheorem{definition}{Definition}}{}
\@ifundefined{lemma}{\newtheorem{lemma}[definition]{Lemma}}{}
\makeatother
\def\e{\varepsilon}
\title[Asymptotic analysis of fractional Sobolev spaces on thin fil]{Asymptotic analysis of fractional Sobolev spaces on thin films in the low-integrability regime}
\author{Andrea Braides, Andrea Pinamonti, Margherita Solci}
\date{Source: arXiv:2512.10620v2 (CC BY 4.0)}
\begin{document}
\maketitle

\noindent\emph{Source and licence.} Asymptotic analysis of fractional Sobolev spaces on thin films in the low-integrability regime. Version arXiv:2512.10620v2 (CC BY 4.0), distributed under the Creative Commons Attribution 4.0 International licence, as recorded in the arXiv licence metadata for this exact version and shown on the abstract page https://arxiv.org/abs/2512.10620v2. Full rights notice: licenses/arxiv-2512.10620-CC-BY-4.0.txt.\par\smallskip

\noindent\textit{Editorial scope.} Counted unit: the Lemma whose source label is \texttt{vert-lemma} in the article (source0.tex lines 344--361, source label \texttt{vert-lemma}), delivered together with the article's own complete original proof of it (source0.tex lines 363--486, one \texttt{proof} environment of 124 lines). No mathematics is written, repaired or re-derived: statement and proof are the article's own text line for line. Required context retained uncounted: none. Every definition, display and label of those lines is retained unchanged. Omitted from this package and not redistributed: 1--343, 362--362, 487--798. Nothing inside a retained excerpt is omitted or altered: every retained body is delivered exactly as the article's lines are written, so no omission receipt of any kind accompanies this package. Citations: the retained excerpts carry no citation command at all, so no bibliography entry of the article is transcribed and no reference is redistributed. Reference adaptation: none was needed and none was made; every reference inside the delivered unit resolves inside it, so no unresolved reference or citation is produced. Printed slips kept literal: no printed word, symbol, hypothesis or number of the counted statement or of its proof was altered. Layout and class: the article's own class is \texttt{article} with packages including graphicx, amsmath, amsthm, amssymb, color, epstopdf, pdfsync, makeidx, geometry; the wrapper is the lane's pinned \texttt{amsart} editorial wrapper at 10pt on A4, which restates the theorem-like environments the retained text needs and adds the article's own macro definitions that the retained text needs (\texttt{\textbackslash e}), with extra wrapper packages (bm, bbm, dsfont, xspace, cancel, mathtools). The delivered numbering is the unit's own. Assets: no retained span contains a figure environment, a graphics-inclusion command, a PDF-page inclusion, a table or a TikZ picture; no image file is copied into this package, the package declares no asset dependency and no image file is redistributed with it. Licence: version arXiv:2512.10620v2 of this article is distributed under the Creative Commons Attribution 4.0 International licence, as recorded in the arXiv licence metadata for this exact version and shown on the abstract page https://arxiv.org/abs/2512.10620v2. A full rights notice is delivered at licenses/arxiv-2512.10620-CC-BY-4.0.txt. Characters: every retained excerpt is pure ASCII, so the delivered engine, XeLaTeX with its default Latin Modern text font, reports no missing character.
\begin{lemma}[estimate of the seminorm in the thin direction]\label{vert-lemma}
Let $u_\e\colon\Omega_\e\to \mathbb R$ be such that there exists $L>0$ such that for all $\e>0$ 
\begin{equation}\label{equilip}
|u_\e(x^\prime, x_d)-u_\e(y^\prime, x_d)|\leq L|x^\prime-y^\prime|
\end{equation}
for all $x^\prime,y^\prime\in\omega$ and $x_d\in (0,\e)$. 
Then, with fixed $\eta>0$ and $\tau>0$,  
\begin{eqnarray}\label{vert-est1}
&&(1+\eta)\lfloor u_\e\rfloor_{s_\e}^2(\Omega_\e)+\frac{\e\phi(\e)}{\eta}\geq C_{s_\e,d} \int_{\omega_\tau}\int_0^\e\int_0^\e\frac{|u_\e(x^\prime, x_d)-u_\e(x^\prime, y_d)|^2}{|x_d-y_d|^{1+2s_\e}}\, dx_d\, dy_d\, dx^\prime\qquad \\
&&\lfloor u_\e\rfloor_{s_\e}^2(\Omega^\tau_\e)-\frac{\e\phi(\e)}{\eta}\leq (1+\eta)C_{s_\e,d} \int_{\omega_\tau}\int_0^\e\int_0^\e\frac{|u_\e(x^\prime, x_d)-u_\e(x^\prime, y_d)|^2}{|x_d-y_d|^{1+2s_\e}}\, dx_d\, dy_d\, dx^\prime,\qquad \label{vert-est2}
\end{eqnarray}
where 
\begin{equation}\label{defC}
C_{s,d}=\int_{\mathbb R^{d-1}}
\frac{1}{(1+|\xi|^2)^{\frac{d}{2}+s}}\, d\xi 
\end{equation}
and $\phi$ only depends on $L$, $\omega$, $\tau$, and on the supremum of the $\infty$ norms of $u_\e$, and it is such that $\lim\limits_{\e\to 0}\phi(\e)=0$. 
\end{lemma}

\begin{proof} 
We extend $u_\e$ to $\mathbb R^{d-1}\times(0,\e)$ by setting $u_\e(x^\prime,x_d)=0$ for $x^\prime\in\mathbb R^{d-1}\setminus\omega$. 
Note that, by a change of variable, 
for all $x_d\neq y_d$ and $x^\prime\in\omega$ 
we have 
\begin{eqnarray}
\frac{C_{s,d}}{|x_d-y_d|^{1+2s}}&=&
\int_{\mathbb R^{d-1}} \frac{|x_d-y_d|^{d-1}}{(|x_d-y_d|^2+(|x_d-y_d||\xi|)^2)^{\frac{d}{2}+s}}\, d\xi\nonumber\\
&=&
\int_{\mathbb R^{d-1}} \frac{1}{(|x_d-y_d|^2+|x^\prime-y^\prime|^2)^{\frac{d}{2}+s}}\, dy^\prime,\label{rd1}
\end{eqnarray} 
so that, by using Fubini's Theorem, we can write  
\begin{eqnarray}
&&\hspace{-2cm}C_{s_\e,d}\int_{\omega_\tau}\int_0^\e\int_0^\e\frac{|u_\e(x^\prime, x_d)-u_\e(x^\prime, y_d)|^2}{|x_d-y_d|^{1+2s_\e}}\, cx_d\, dy_d\, dx^\prime\nonumber\\
&=&
\int_{\omega_\tau}\int_0^\e\int_0^\e\int_{\mathbb R^{d-1}} \frac{|u_\e(x^\prime, x_d)-u_\e(x^\prime, y_d)|^2}{(|x_d-y_d|^2+|x^\prime-y^\prime|^2)^{\frac{d}{2}+s_\e}}\, dy^\prime\, dx_d\, dy_d\, dx^\prime.\label{cvar}
\end{eqnarray}
Setting, for $A$ measurable subset of $\mathbb R^{d-1}$, 
\begin{eqnarray*}
&&I^{\rm v}_{\e,\tau}(A)=\int_{\omega_\tau}\int_0^\e\int_0^\e\int_A \frac{|u_\e(x^\prime, x_d)-u_\e(x^\prime, y_d)|^2}{(|x_d-y_d|^2+|x^\prime-y^\prime|^2)^{\frac{d}{2}+s_\e}}\, dy^\prime\, dx_d\, dy_d\, dx^\prime\\ 
&&I^{\rm h}_{\e,\tau}(A)=\int_{\omega_\tau}\int_0^\e\int_0^\e\int_A \frac{|u_\e(y^\prime, y_d)-u_\e(x^\prime, y_d)|^2}{(|x_d-y_d|^2+|x^\prime-y^\prime|^2)^{\frac{d}{2}+s_\e}}\, dy^\prime\, dx_d\, dy_d\, dx^\prime  \\
&&I_{\e,\tau}(A)= \int_{\omega_\tau}\int_0^\e\int_0^\e\int_A\frac{|u_\e(x^\prime, x_d)-u_\e(y^\prime, y_d)|^2}{(|x_d-y_d|^2+|x^\prime-y^\prime|^2)^{\frac{d}{2}+s_\e}}\, dy^\prime\, dx_d\, dy_d\, dx^\prime, 
\end{eqnarray*}
a triangular inequality gives 
\begin{eqnarray}
&&I_{\e,\tau}^{\rm v}(\mathbb R^{d-1})\leq (1+\eta)I_{\e,\tau}(\mathbb R^{d-1})+\frac{1}{\eta}I_{\e,\tau}^{\rm h}(\mathbb R^{d-1})\label{tr1}\\
&&I_{\e,\tau}(\mathbb R^{d-1})\leq (1+\eta)I_{\e,\tau}^{\rm v}(\mathbb R^{d-1})+\frac{1}{\eta}I_{\e,\tau}^{\rm h}(\mathbb R^{d-1})\label{tr2}
\end{eqnarray}
for all $\eta>0$. 
We now show that 
\begin{equation}\label{stima-or}
I_{\e,\tau}^{\rm h}(\mathbb R^{d-1})\leq\e\phi(\e),
\end{equation}
where $\phi$ depends on $L$, $\omega$ and $\tau$, and $\lim\limits_{\e\to 0}\phi(\e)=0$.
Noting that for all $x^\prime\in\omega_\tau$ 
\begin{equation*}
\int_{\mathbb R^{d-1}\setminus\omega} \frac{1}{(|x_d-y_d|^2+|x^\prime-y^\prime|^2)^{\frac{d}{2}+s_\e}} dy^\prime
%\leq \int_{\mathbb R^{d-1}\setminus\omega} \frac{1}{|x^\prime-y^\prime|^{d+2s_\e}} dy^\prime
\leq\int_{\mathbb R^{d-1}\setminus B_\tau(0)} \frac{1}{|\xi^\prime|^{d+2s_\e}}\, d\xi^\prime
=\frac{\sigma_{d-1}}{(1+2s_\e)\tau^{1+2s_\e}}, 
\end{equation*}
where $\sigma_{d-1}=\mathcal H^{d-2}(\partial B_1(0))$, 
we get 
\begin{equation}
I^{\rm h}_{\e,\tau}(\mathbb R^{d-1}\setminus \omega)
\leq \frac{\sigma_{d-1}\|u_\e\|^2_\infty |\omega_\tau|}{\tau^{1+2s_\e}}\e^2\label{st1}\\
%&&\int_{\omega_\tau}\int_0^\e\int_0^\e\int_{\mathbb R^{d-1}\setminus\omega}\frac{|u_\e(y^\prime, y_d)-u_\e(x^\prime, y_d)|^2}{(|x_d-y_d|^2+|x^\prime-y^\prime|^2)^{\frac{d}{2}+s}}\, dy^\prime\, dx_d\, dy_d\, dx^\prime
%\leq \frac{\sigma_{d-1}\|u_\e\|^2_\infty |\omega_\tau|}{\tau^{1+2s_\e}}\e^2\label{st2}. 
\end{equation}
Hence, we only have to estimate the integral in the set where $y^\prime\in\omega$. 
We have 
\begin{eqnarray*}
I^{\rm h}_{\e,\tau}(\omega)
&\leq& 
2L^2\int_{\omega_\tau}\int_0^\e\int_{\omega}|x^\prime-y^\prime|^2\int_0^\e\frac{1}{(t^2+|x^\prime-y^\prime|^2)^{\frac{d}{2}+s}}\, dt\, dy^\prime\, dy_d\, dx^\prime\\
&\leq& 
2L^2\e\phi(\e)
%\int_{\omega}\int_{\omega}|x^\prime-y^\prime|^2\psi_\e(|x^\prime-y^\prime|)\, dy^\prime\, dx^\prime
\end{eqnarray*}
where 
\begin{equation}\label{defphi}\phi(\e)=
\int_{\omega}\int_{\omega}|x^\prime-y^\prime|^2
\int_0^\e\frac{1}{(t^2+|x^\prime-y^\prime|^2)^{\frac{d}{2}+s}}\, dt
\, dy^\prime\, dx^\prime.
\end{equation}
Since the integrand in \eqref{defphi} pointwise converges to $0$ almost everywhere, and it is estimated by 
\begin{eqnarray*}
%\int_0^\e \frac{1}{(|x_d-y_d|^2+|x^\prime-y^\prime|^2)^{\frac{d}{2}+s_\e}}\leq 2
|x^\prime-y^\prime|^2\int_0^\e\frac{1}{(t^2+|x^\prime-y^\prime|^2)^{\frac{d}{2}+s}}\, dt
%&=& 2\frac{|x^\prime-y^\prime|}{|x^\prime-y^\prime|^{d+2s}}
%\int_0^{+\infty}\frac{1}{(1+w^2)^{\frac{d}{2}+s}} d w
&\leq&|x^\prime-y^\prime|^{3-d-2s_\e} \int_0^{+\infty}\frac{1}{(1+w^2)^{\frac{d}{2}+s_\e}} d w \\
%\int_0^{+\infty}\frac{1}{(1+w^2)^{\frac{d}{2}+s}} d w
&\leq&\hbox{\rm diam}(\omega)^{1-2s_\e} |x^\prime-y^\prime|^{2-d}\int_0^{+\infty}\frac{1}{(1+w^2)^{\frac{d}{2}+s_\e}} d w, 
\end{eqnarray*}
with
$$\int_{\omega}\int_{\omega}|x^\prime-y^\prime|^{2-d}\, dx^\prime\, dy^\prime<+\infty,$$ 
by Lebesgue's Theorem we obtain 
$$\lim_{\e\to 0}\phi(\e)=0.$$
%\begin{eqnarray*}
%C\int_{\omega_\tau}\int_0^\e\int_{\mathbb R^{d-1}\setminus\omega}
%\frac{|u_\e(x^\prime, y_d)|^2}{|x^\prime-y^\prime|^{d-1+2s}}\, dy^\prime\, dy_d\, %dx^\prime&\leq& C\|u_\e\|_{\infty}^2|\omega|\e \int_{\mathbb R^{d-1}\setminus %B_\tau(0)}\frac{1}{|\xi|^{d-1+2s}}\, d\xi
%\end{eqnarray*}
Since \eqref{st1} holds, this implies \eqref{stima-or}.

Noting that \eqref{st1} holds also for $I_{\e,\tau}(\mathbb R^{d-1}\setminus\omega)$, we can write
%\begin{eqnarray}\label{stima-semi}
%&&\hspace{-2cm}\int_{\omega_\tau}\int_0^\e\int_0^\e\int_{\mathbb R^{d-1}} \frac{|u_\e(x^\prime, x_d)-u_\e(y^\prime, y_d)|^2}{(|x_d-y_d|^2+|x^\prime-y^\prime|^2)^{\frac{d}{2}+s}}\, dy^\prime\, dx_d\, dy_d\, dx^\prime\nonumber\\
%&=& \int_{\omega_\tau}\int_0^\e\int_0^\e\int_{\omega} \frac{|u_\e(x^\prime, x_d)-%u_\e(y^\prime, y_d)|^2}{(|x_d-y_d|^2+|x^\prime-y^\prime|^2)^{\frac{d}{2}+s}}\, %dy^\prime\, dx_d\, dy_d\, dx^\prime+r_\e, 
%\nonumber\\
%&&+\frac{C}{\tau^{1+2s_\e}}\int_{\omega_\tau}\int_0^\e\int_0^\e |u_\e(x^\prime, x_d)|%^2\, dx_d\, dy_d\, dx^\prime\nonumber\\
%&\leq& \lfloor u_\e\rfloor_{s_\e}^2(\Omega_\e) + C\e^2\|u_\e\|^2_\infty |\omega|%\tau^{-1-2s_\e}, 
%\end{eqnarray}
\begin{equation}\label{stima-semi}
I_{\e,\tau}(\mathbb R^{d-1})
= I_{\e,\tau}(\omega)+r_\e, 
%\nonumber\\
%&&+\frac{C}{\tau^{1+2s_\e}}\int_{\omega_\tau}\int_0^\e\int_0^\e |u_\e(x^\prime, x_d)|^2\, dx_d\, dy_d\, dx^\prime\nonumber\\
%&\leq& \lfloor u_\e\rfloor_{s_\e}^2(\Omega_\e) + C\e^2\|u_\e\|^2_\infty |\omega|\tau^{-1-2s_\e}, 
\end{equation}
with $|r_\e|\leq \sigma_{d-1}\|u_\e\|^2_\infty |\omega_\tau|\tau^{-1-2s_\e}\e^2$. 

%By using Fubini's Theorem and \eqref{rd1}, we can then estimate 
%\begin{eqnarray*}
%&&\hspace{-2cm}
%C_{s_\e,d}\int_{\omega_\tau}\int_0^\e\int_0^\e\frac{|u_\e(x^\prime, x_d)-u_\e(x^\prime, %y_d)|^2}{|x_d-y_d|^{1+2s_\e}}\, cx_d\, dy_d\, dx^\prime\\
%&=&
%\int_{\omega_\tau}\int_0^\e\int_0^\e\int_{\mathbb R^{d-1}} \frac{|u_\e(x^\prime, x_d)-u_\e(x^\prime, y_d)|^2}{(|x_d-y_d|^2+|x^\prime-y^\prime|^2)^{\frac{d}{2}+s_\e}}\, dy^\prime\, dx_d\, dy_d\, dx^\prime\\
%&\leq&(1+\eta)\int_{\omega_\tau}\int_0^\e\int_0^\e\int_{\mathbb R^{d-1}} \frac{|u_\e(x^\prime, x_d)-u_\e(y^\prime, y_d)|^2}{(|x_d-y_d|^2+|x^\prime-y^\prime|^2)^{\frac{d}{2}+s_\e}}\, dy^\prime\, dx_d\, dy_d\, dx^\prime\\
%&&+\frac{1}{\eta}
%\int_{\omega_\tau}\int_0^\e\int_0^\e\int_{\mathbb R^{d-1}} \frac{|u_\e(y^\prime, y_d)-%u_\e(x^\prime, y_d)|^2}{(|x_d-y_d|^2+|x^\prime-y^\prime|^2)^{\frac{d}{2}+s_\e}}\, %dy^\prime\, dx_d\, dy_d\, dx^\prime.
%\end{eqnarray*}
Recalling \eqref{tr1} and \eqref{tr2}, estimates \eqref{stima-or} and \eqref{stima-semi} imply that 
\begin{eqnarray*}
I_{\e,\tau}^{\rm v}(\mathbb R^{d-1})&\leq&(1+\eta)(I_{\e,\tau}(\omega)+|r_\e|)+\frac{\e\phi(\e)}{\eta}\\
&\leq&(1+\eta)\lfloor u_\e\rfloor_{s_\e}^2(\Omega_\e)
+C\e^2\tau^{-1-2s_\e}
+\frac{\e\phi(\e)}{\eta}
\end{eqnarray*}
\begin{eqnarray*}
\lfloor u_\e\rfloor_{s_\e}^2(\Omega_\e^\tau)-C\e^2\tau^{-1-2s_\e}&\leq& I_{\e,\tau}(\omega)-|r_\e|\\
&\leq&(1+\eta)I_{\e,\tau}^{\rm v}(\mathbb R^{d-1})+\frac{\e\phi(\e)}{\eta}, 
\end{eqnarray*}
and since \eqref{cvar} holds the claim follows. \end{proof}

\end{document}
